\section{语言学为什么仍然重要}

\subsection{语言学组织 NLP 的未来}

这节课有一个很明确的立场：语言学不是过时的背景知识，而是理解和评估 NLP 的关键工具。即使具体的语法理论流派未必是今天最需要纠结的点，语言学提供的那些宽泛区分依然很有用。

\begin{figure}[H]
\centering
\includegraphics[width=0.95\textwidth]{slides-images/slide-030.jpg}
\caption{语言学不是 NLP 的装饰，而是组织问题、分析系统、制定评测标准的工具\protect\footnotemark}
\end{figure}
\footnotetext{Slides 强调：结构、篇章、推理、语用等语言学概念在 2020s NLP 中仍然有用。}

语言学在今天的重要性，主要体现在三个层面：
\begin{enumerate}
    \item \textbf{定义任务}：我们到底要让系统理解什么；
    \item \textbf{分析行为}：模型为什么会在某些现象上失败；
    \item \textbf{设计评测}：不能只看一个总分，必须细分语言现象。
\end{enumerate}

具体来说，现代 NLP 仍然会遇到大量语言学现象：句法结构、篇章结构、NLI、夸张、变体语言、语体、隐喻、指代、预设、语气、风格、语音协同发音等。它们不是“边角料”，而是自然语言本身的一部分。

\begin{knowledgebox}{语言学不是只有句法}
如果只盯着“句子树”或“依存边”，就会误以为语言学只关心形式结构。实际上，语言学还关心语义、语用、篇章、话语角色、说话方式、社会语境和语音细节。很多模型看似能做 translation 或 summarization，但在这些层次上仍然脆弱。
\end{knowledgebox}

\subsection{本章小结}

语言学并没有被神经模型取代，而是在新的 NLP 时代变成了分析框架和问题分类器。它帮助我们把“模型错了”这件事说得更具体。

%% ========== 符号系统与神经系统 ==========

